\documentclass[../../main.tex]{subfiles}

\begin{document}

\chapter{Segmentasi Memori dan Alamat Fisik}
\label{chap:segmentasi}

% =========================================================
% OUTCOME PEMBELAJARAN KHUSUS BAB
% =========================================================

\begin{subcpmkbox}
\SubCPMK{1.3} --- Menerapkan konsep segmentasi memori: menghitung
alamat fisik 20-bit dari pasangan \eng{segment:offset}, menjelaskan
paragraf 16-byte, dan menganalisis skema pengalamatan 1 MB.

\vspace{0.4em}
\textbf{CPMK yang didukung:} \CPMK{01} \quad\textbullet\quad
\textbf{CPL yang didukung:} \CPL{02}
\end{subcpmkbox}

\begin{tujuanbox}
Setelah mempelajari bab ini, mahasiswa mampu:
\begin{enumerate}
    \item menjelaskan mengapa prosesor 8086 yang berbus alamat 20 bit
          tetap memakai register segmen 16 bit;
    \item membedakan alamat logis (\eng{segment:offset}) dan alamat
          fisik, serta menuliskan keduanya dengan notasi yang benar;
    \item menghitung alamat fisik 20-bit dari pasangan
          \eng{segment:offset} memakai rumus
          $(\text{segmen} \times \hex{10}) + \text{offset}$;
    \item menjelaskan makna paragraf 16-byte dan hubungannya dengan
          ukuran segmen 64 KB;
    \item menghitung banyaknya pasangan \eng{segment:offset} yang
          menunjuk satu alamat fisik yang sama
          (\eng{overlapping segment});
    \item menyebutkan pembagian peta memori 1 MB beserta kegunaan
          setiap area (vektor interupsi, memori konvensional, area
          video, dan ROM BIOS);
    \item menganalisis kasus batas perhitungan alamat yang melewati
          batas 1 MB (\eng{wrap-around});
    \item memverifikasi hasil perhitungan alamat dengan
          \eng{debugger} Emu8086 atau DEBUG.
\end{enumerate}
\end{tujuanbox}

\begin{materibox}
\begin{enumerate}
    \item Masalah yang melahirkan segmentasi: register 16-bit
          (maksimum 64 KB) berhadapan dengan bus alamat 20-bit
          (1 MB).
    \item Alamat logis (\eng{segment:offset}) dan alamat fisik;
          notasi penulisan \code{1234h:5678h}.
    \item Rumus alamat fisik
          $(\text{segmen} \times \hex{10}) + \text{offset}$ dan
          alasan penggunaan pengali 16.
    \item Paragraf 16-byte dan struktur segmen 64 KB.
    \item \eng{Overlapping segment}: banyak pasangan menunjuk satu
          alamat fisik; cara menghitung jumlah kemungkinannya.
    \item Peta memori 1 MB: tabel vektor interupsi, area data BIOS,
          memori konvensional 640 KB, area video, dan ROM BIOS.
    \item Kasus batas: alamat yang melewati batas 1 MB dan
          \eng{wrap-around}.
    \item Verifikasi alamat dengan DEBUG dan Emu8086
          (\lstasm|D segmen:offset|).
\end{enumerate}
\end{materibox}

\subfile{section-01}
\subfile{section-02}
\subfile{section-03}

\section{Contoh Program}
\label{sec:contoh-program}

Dua contoh berikut memperlihatkan bagaimana nilai segmen dan offset
diperoleh dari dalam program, lalu bagaimana keduanya dipakai untuk
menghitung alamat fisik. Keduanya memakai prosedur bantu
\code{cetak\_hex} yang mencetak isi register \reg{AX} sebagai empat
digit heksadesimal, sehingga mahasiswa dapat membaca nilai segmen
dan offset yang sesungguhnya tanpa harus membuka emulator.

\subsection{Contoh 1: Menampilkan Offset Variabel}

Program pada Kode~\ref{lst:alamat-fisik} mendeklarasikan dua
variabel pada segmen data, lalu menampilkan offset keduanya dengan
operator \code{OFFSET}.

\lstinputlisting[
    language=MASM,
    caption={Program \file{alamat-fisik.asm} menampilkan offset
    variabel pada segmen data.},
    label={lst:alamat-fisik}
]{\CodeDir/alamat-fisik.asm}

Bagian yang perlu diperhatikan dari program tersebut adalah cara
nilai offset diperoleh:

\begin{lstlisting}[language=MASM, numbers=none]
MOV AX, OFFSET var_byte    ; AX = offset var_byte
MOV AX, OFFSET var_word    ; AX = offset var_word
\end{lstlisting}

Operator \code{OFFSET} menghasilkan jarak variabel tersebut dari
awal segmen data, bukan alamat fisiknya. Nilai inilah yang menjadi
bagian kedua pada pasangan \eng{segment:offset}. Karena variabel
\code{var\_byte} dideklarasikan lebih dahulu, offsetnya nol,
sedangkan \code{var\_word} berada tepat sesudahnya dengan offset
satu. Keluaran program tampak pada Kode~\ref{lst:keluaran-offset}.

\begin{session}[caption={Keluaran program \file{alamat-fisik.asm}.},
label={lst:keluaran-offset}]
Offset var_byte = 0000
Offset var_word = 0001
\end{session}

\begin{infobox}
Nilai offset yang tampil dapat berbeda satu atau dua byte bila
assembler yang dipakai menyelaraskan data ke alamat genap. Yang
penting bukan angkanya, melainkan caranya membaca: angka yang tampil
adalah offset, dan nomor segmennya harus dibaca terpisah dari
register \reg{DS}.
\end{infobox}

Untuk memperoleh alamat fisik, mahasiswa harus mengetahui nilai
\reg{DS} saat program berjalan. Misalkan pada Emu8086 atau DEBUG
terbaca \reg{DS} $= \hex{076A}$, maka kedua alamat fisik tersebut
dapat dihitung seperti pada Tabel~\ref{tab:analisis-offset}.

\begin{table}[H]
    \centering
    \caption{Perhitungan alamat fisik variabel pada program pertama
    dengan asumsi \reg{DS} $= \hex{076A}$.}
    \label{tab:analisis-offset}
    \small
    \begin{tabular}{@{}p{2.3cm}p{1.8cm}p{1.5cm}p{4.1cm}p{1.9cm}@{}}
        \toprule
        \textbf{Variabel} & \textbf{Offset} & \textbf{DS} &
        \textbf{Perhitungan} & \textbf{Alamat fisik} \\
        \midrule
        \code{var\_byte} & \hex{0000} & \hex{076A} &
        \hex{076A0} + \hex{0000} & \hex{076A0} \\
        \addlinespace
        \code{var\_word} & \hex{0001} & \hex{076A} &
        \hex{076A0} + \hex{0001} & \hex{076A1} \\
        \bottomrule
    \end{tabular}
\end{table}

\subsection{Contoh 2: Membaca Nilai Segmen dan Menghitung Alamat}

Program kedua, yang ditampilkan pada Kode~\ref{lst:hitung-alamat},
melangkah lebih jauh: ia membaca nomor segmen dengan operator
\code{seg}, membaca offset, lalu menghitung alamat fisiknya di dalam
program.

\lstinputlisting[
    language=MASM,
    caption={Program \file{hitung-alamat.asm} menghitung alamat
    fisik 20-bit dari pasangan segmen:offset.},
    label={lst:hitung-alamat}
]{\CodeDir/hitung-alamat.asm}

Inti perhitungannya terletak pada empat instruksi berikut.

\begin{lstlisting}[language=MASM, numbers=none]
MOV AX, seg var_nilai      ; AX = nomor segmen var_nilai
MOV DX, AX                 ; DX akan menerima bit 16..19
MOV CL, 4
SHR DX, CL                 ; DX = segmen / 16
SHL AX, CL                 ; AX = segmen x 16
ADD AX, BX                 ; AX = (segmen x 16) + offset
ADC DX, 0                  ; tambahkan carry ke DX
\end{lstlisting}

Perhatikan bahwa pergeseran ke kiri empat bit (\lstasm|SHL AX, CL|)
adalah bentuk lain dari perkalian dengan 16, dan pergeseran ke
kanan empat bit (\lstasm|SHR DX, CL|) memindahkan bagian segmen yang
melewati 16 bit ke register \reg{DX}. Karena hasil penjumlahan
dapat mencapai 20 bit, alamat fisik disimpan pada pasangan register
\reg{DX}:\reg{AX}, dengan \reg{DX} memuat bit 16 sampai 19 dan
\reg{AX} memuat bit 0 sampai 15. Keluaran program tampak pada
Kode~\ref{lst:keluaran-hitung}.

\begin{session}[caption={Keluaran program \file{hitung-alamat.asm}.},
label={lst:keluaran-hitung}]
Segmen: 076A
Offset: 0000
Alamat fisik: 000076A0
\end{session}

Keluaran tersebut menunjukkan tiga nilai sekaligus: nomor segmen
\hex{076A}, offset \hex{0000}, dan alamat fisik \hex{076A0} yang
ditulis delapan digit karena menempati pasangan \reg{DX}:\reg{AX}.
Mahasiswa dapat mengubah isi \reg{AX} sebelum perhitungan, memindahkan
variabel \code{var\_nilai} ke posisi lain pada segmen data, atau
mengganti segmen yang dipakai, lalu memeriksa apakah alamat fisik
yang dihasilkan tetap sesuai dengan perhitungan manual.

\begin{infobox}
Operator \code{seg} hanya menghasilkan \emph{nomor} segmen, bukan
alamat dasarnya. Pada contoh di atas, \lstasm|MOV AX, seg var_nilai|
memberi \hex{076A}, sedangkan alamat dasar segmen tersebut adalah
\hex{076A0}. Perbedaan satu digit nol inilah yang menjadi sumber
kesalahan paling sering pada perhitungan alamat.
\end{infobox}

\section{Praktikum}
\label{sec:praktikum-bab6}

\begin{praktikumbox}
\vspace{1.1em}
\textbf{Judul:} Menghitung dan memverifikasi alamat fisik dengan
\eng{debugger}.\\[0.3em]
\textbf{Sub-CPMK:} 1.4 \quad\textbullet\quad
\textbf{Alokasi waktu:} 100 menit \quad\textbullet\quad
\textbf{Bentuk:} individu
\end{praktikumbox}

\subsection{Tujuan Praktikum}

Setelah menyelesaikan praktikum ini, mahasiswa mampu menghitung
alamat fisik dari pasangan \eng{segment:offset}, membuktikan hasilnya
dengan DEBUG dan emu8086, serta menunjukkan \eng{overlapping} segmen.

\subsection{Alat dan Bahan}

\begin{itemize}
    \item Komputer dengan emu8086 terpasang.
    \item DOSBox dengan TASM, TLINK, dan DEBUG.
    \item Berkas \file{alamat-fisik.asm} dan \file{hitung-alamat.asm}
          dari \file{code/asm/}.
    \item Kalkulator heksadesimal atau lembar kerja perhitungan.
\end{itemize}

\subsection{Langkah Kerja}

\begin{enumerate}
    \item \textbf{Merakit dan menjalankan program pertama.}
    \begin{enumerate}
        \item Buka DOSBox, lalu masuk ke direktori kerja.
        \item Rakit \file{alamat-fisik.asm} dengan
              \file{tasm alamat-fisik.asm}, lalu link dengan
              \file{tlink alamat-fisik.obj}.
        \item Jalankan program dan catat kedua nilai offset yang
              tampil.
    \end{enumerate}

    \item \textbf{Menghitung alamat fisik secara manual.}
    \begin{enumerate}
        \item Jalankan \file{debug alamat-fisik.exe}.
        \item Ketik \file{r}, lalu catat nilai \reg{DS}.
        \item Hitung alamat fisik setiap variabel dengan rumus
              $(\text{segmen} \times \hex{10}) + \text{offset}$.
              Tuliskan langkahnya, bukan hanya hasil akhir.
        \item Susun hasil pada tabel seperti
              Tabel~\ref{tab:analisis-offset}.
    \end{enumerate}

    \item \textbf{Memverifikasi dengan DEBUG.}
    \begin{enumerate}
        \item Tampilkan isi memori dengan
              \lstasm|D DS:0000 L 08|.
        \item Bandingkan isi yang tampil dengan nilai variabel pada
              program: \code{var\_byte} bernilai \hex{10} dan
              \code{var\_word} bernilai \hex{1234}.
        \item Perhatikan urutan byte pada \code{var\_word}. Jelaskan
              mengapa byte rendah tampil lebih dahulu.
    \end{enumerate}

    \item \textbf{Membuktikan overlapping.}
    \begin{enumerate}
        \item Bila \reg{DS} $= \hex{076A}$ dan offset
              \code{var\_word} $= \hex{0001}$, alamat fisiknya adalah
              \hex{076A1}.
        \item Tuliskan pasangan lain yang menunjuk alamat tersebut,
              misalnya \code{0769h:0011h}.
        \item Tampilkan kedua pasangan dengan perintah
              \lstasm|D| dan pastikan isinya sama.
        \item Jelaskan mengapa dua pasangan berbeda dapat menunjuk
              lokasi yang sama.
    \end{enumerate}

    \item \textbf{Menjalankan program kedua.}
    \begin{enumerate}
        \item Rakit dan jalankan \file{hitung-alamat.asm}.
        \item Bandingkan alamat fisik yang dihitung program dengan
              hasil perhitungan manual.
        \item Tambahkan satu variabel baru di depan \code{var\_nilai},
              rakit ulang, lalu jelaskan perubahan offset yang
              terjadi.
    \end{enumerate}

    \item \textbf{Kuis 1.}
    \begin{enumerate}
        \item Kerjakan kuis secara individu pada akhir pertemuan.
        \item Materi kuis mencakup perhitungan alamat fisik,
              paragraf 16 byte, \eng{overlapping} segmen, dan
              pembagian memori 1 MB.
    \end{enumerate}
\end{enumerate}

\subsection{Hasil yang Dikumpulkan}

\begin{enumerate}
    \item Lembar perhitungan manual alamat fisik untuk minimal enam
          pasangan \eng{segment:offset}.
    \item Tangkapan layar keluaran \file{alamat-fisik.asm} dan
          \file{hitung-alamat.asm}.
    \item Tangkapan layar sesi DEBUG dengan perintah \file{r} dan
          \file{d}.
    \item Bukti pengamatan \eng{overlapping} beserta penjelasan.
    \item Jawaban Kuis 1.
\end{enumerate}

\clearpage
\begin{asesmenbox}
\vspace{1.1em}
\textbf{Asesmen praktikum 4 (bobot 5\% dari komponen praktikum)} dan
\textbf{Kuis 1 (bobot 10\% dari komponen kuis)}.

\begin{tabular}{@{}p{6.4cm}cp{4.4cm}@{}}
    \toprule
    \textbf{Indikator} & \textbf{Bobot} & \textbf{Kriteria} \\
    \midrule
    Ketepatan perhitungan alamat fisik & 35\% &
    Seluruh langkah ditulis dan hasilnya benar \\
    \addlinespace
    Ketepatan verifikasi dengan DEBUG & 30\% &
    Perintah yang dipakai tepat dan isinya ditafsirkan dengan benar \\
    \addlinespace
    Analisis \eng{overlapping} & 20\% &
    Penjelasan pasangan alamat dan alamat fisik tepat \\
    \addlinespace
    Kerapian dokumentasi & 15\% &
    Langkah kerja terdokumentasi dan dapat diulang \\
    \bottomrule
\end{tabular}
\end{asesmenbox}

\section{Latihan}
\label{sec:latihan-bab6}

\begin{enumerate}
    \item \textbf{Konseptual.} Jelaskan mengapa 8086 yang berbus
          alamat 20 bit memakai register segmen 16 bit. Apa yang
          akan hilang bila prosesor hanya memakai satu register
          16 bit untuk menunjuk alamat?

    \item \textbf{Perhitungan.} Hitung alamat fisik dari pasangan
          berikut. Tuliskan langkahnya.
          \begin{enumerate}
              \item \code{0000h:1A00h}
              \item \code{1500h:0100h}
              \item \code{9FFFh:000Fh}
              \item \code{B800h:0060h}
              \item \code{ABCDh:1234h}
          \end{enumerate}

    \item \textbf{Overlapping.} Alamat fisik \hex{00080} dapat
          dinyatakan oleh beberapa pasangan \eng{segment:offset}.
          \begin{enumerate}
              \item Tentukan banyaknya pasangan dengan pola
                    $\lfloor F/16 \rfloor + 1$.
              \item Tuliskan seluruh pasangan tersebut.
              \item Jelaskan mengapa jumlah pasangan untuk
                    \hex{00080} lebih banyak daripada \hex{00050}.
          \end{enumerate}

    \item \textbf{Peta memori.} Berdasarkan
          Tabel~\ref{tab:peta-memori}, jawab pertanyaan berikut.
          \begin{enumerate}
              \item Wilayah mana yang disebut memori konvensional,
                    dan berapa ukurannya?
              \item Apa isi area \hex{A0000}--\hex{BFFFF}?
              \item Di wilayah mana ROM BIOS berada?
          \end{enumerate}

    \item \textbf{Penulisan program.} Tulislah program yang
          mendeklarasikan tiga variabel pada segmen data, lalu
          menampilkan offset ketiganya dengan operator
          \code{OFFSET}. Gunakan prosedur \code{cetak\_hex} pada
          Kode~\ref{lst:alamat-fisik} sebagai acuan.

    \item \textbf{Penulisan program.} Tulislah program yang
          menghitung alamat fisik pasangan \code{1234h:5678h}
          secara langsung, yaitu memuat kedua nilai sebagai
          konstanta ke register, lalu melakukan pergeseran dan
          penjumlahan seperti pada Kode~\ref{lst:hitung-alamat}.
          Program harus mencetak \hex{179B8}.

    \item \textbf{Analisis kesalahan.} Seorang mahasiswa menghitung
          alamat fisik \code{1234h:5678h} dan memperoleh
          \hex{68AC}. Ia juga melihat DEBUG menampilkan
          \code{076A:0000} dan menyimpulkan alamat fisiknya
          \hex{0000}.
          \begin{enumerate}
              \item Tunjukkan kesalahan pada perhitungan pertama
                    dan hitung hasil yang benar.
              \item Jelaskan kesalahan penalaran pada kesimpulan
                    kedua, lalu tentukan alamat fisik yang benar
                    bila \reg{DS} $= \hex{076A}$.
              \item Sebutkan satu cara memeriksa kebenaran hasil.
          \end{enumerate}

    \item \textbf{Konseptual.} Jelaskan hubungan antara 1 MB,
          segmen 64 KB, dan paragraf 16 byte. Mengapa satu alamat
          fisik dapat memiliki paling banyak 4.096 pasangan
          \eng{segment:offset}?
\end{enumerate}

\begin{refleksibox}
\vspace{1.1em}
Jawab pertanyaan berikut pada jurnal praktikum Anda.
\begin{enumerate}
    \item Bagian mana yang paling sulit: aritmetika heksadesimal,
          istilah segmentasi, atau pembacaan peta memori? Mengapa?
    \item Setelah membuktikan \eng{overlapping} dengan DEBUG,
          bagaimana pandangan Anda tentang hubungan alamat logis
          dan alamat fisik?
    \item Konsep apa yang perlu Anda ulangi sebelum mempelajari
          memulai assembly pada bab berikutnya?
\end{enumerate}
\end{refleksibox}

\section{Rangkuman}
\label{sec:rangkuman-bab6}

Pada bab ini telah dibahas:

\begin{itemize}
    \item Register 16 bit hanya menjangkau 64 KB, sedangkan bus
          alamat 8086 selebar 20 bit menjangkau 1 MB. Segmentasi
          menjembatani kedua sisi tersebut.

    \item Alamat logis ditulis sebagai pasangan
          \eng{segment:offset}, misalnya \code{1234h:5678h}.
          Alamat fisik adalah alamat 20 bit yang dikirim ke bus
          alamat.

    \item Alamat fisik dihitung dengan rumus
          $(\text{segmen} \times \hex{10}) + \text{offset}$.
          Pengali \hex{10} sama dengan 16 desimal.

    \item Memori dibagi menjadi paragraf berukuran 16 byte. Nilai
          segmen menyatakan nomor paragraf, sehingga alamat dasar
          segmen selalu kelipatan 16.

    \item Satu segmen menjangkau 64 KB, dari offset \hex{0000}
          sampai \hex{FFFF}. Enam belas segmen menutupi 1 MB.

    \item \eng{Overlapping} segmen membuat satu alamat fisik dapat
          dinyatakan oleh banyak pasangan \eng{segment:offset}.
          Jumlah pasangan untuk alamat $F$ di bawah \hex{10000}
          adalah $\lfloor F/16 \rfloor + 1$, dan maksimumnya 4.096.

    \item Memori 1 MB terbagi menjadi memori konvensional 640 KB,
          buffer video, ROM adaptor, dan ROM BIOS.

    \item Alamat fisik dapat diverifikasi dengan DEBUG atau
          emu8086. Perintah \lstasm|D segmen:offset| menampilkan
          alamat logis; alamat fisiknya dihitung dengan rumus.
\end{itemize}

\begin{checklistbox}
\vspace{1.1em}
Tandai setiap butir yang sudah Anda kuasai.

\begin{itemize}
    \item[\checkbox] Saya dapat menjelaskan alasan segmentasi pada
          8086.
    \item[\checkbox] Saya dapat membedakan alamat logis dan alamat
          fisik.
    \item[\checkbox] Saya dapat menghitung alamat fisik dari
          \eng{segment:offset}.
    \item[\checkbox] Saya dapat menjelaskan paragraf 16 byte dan
          segmen 64 KB.
    \item[\checkbox] Saya dapat menentukan banyaknya pasangan yang
          menunjuk satu alamat fisik.
    \item[\checkbox] Saya dapat menyebutkan pembagian memori 1 MB.
    \item[\checkbox] Saya dapat memverifikasi alamat dengan DEBUG
          atau emu8086.
\end{itemize}

\vspace{0.4em}
Bila masih ada butir yang belum bertanda, ulangi bagian yang
bersesuaian sebelum melanjutkan ke bab berikutnya.
\end{checklistbox}

\begin{infobox}
\vspace{1.1em}
\textbf{Rujukan bab.} Materi pada bab ini diringkas dari
\cite{gunawan2019}, Bagian 2.1.1, hlm.~13--15, dan \cite{sto2001},
BAB II, hlm.~9--13.
\end{infobox}

\section{Evaluasi}
\label{sec:evaluasi-bab6}

\subsection{Soal Pilihan Ganda}

\begin{enumerate}
    \item Rumus yang benar untuk menghitung alamat fisik dari
          pasangan \eng{segment:offset} adalah \dots
          \begin{enumerate}
              \item $\text{segmen} + \text{offset}$
              \item $(\text{segmen} \times \hex{10}) +
                    \text{offset}$
              \item $(\text{segmen} \times \dec{20}) +
                    \text{offset}$
              \item $(\text{offset} \times \hex{10}) +
                    \text{segmen}$
          \end{enumerate}

    \item Nilai segmen \hex{1234} menunjuk alamat dasar \dots
          \begin{enumerate}
              \item \hex{1234} \item \hex{12340}
              \item \hex{01234} \item \hex{23400}
          \end{enumerate}

    \item Alamat fisik yang ditunjuk \code{B800h:0000h} adalah
          \dots
          \begin{enumerate}
              \item \hex{B800} \item \hex{0B800}
              \item \hex{B8000} \item \hex{BB000}
          \end{enumerate}

    \item Satu paragraf pada skema pengalamatan 8086 berukuran
          \dots
          \begin{enumerate}
              \item 16 byte \item 256 byte
              \item 1 KB \item 64 KB
          \end{enumerate}

    \item Jumlah pasangan \eng{segment:offset} terbanyak yang dapat
          menunjuk satu alamat fisik adalah \dots
          \begin{enumerate}
              \item 16 \item 64
              \item 1.024 \item 4.096
          \end{enumerate}
\end{enumerate}

\subsection{Soal Uraian}

\begin{enumerate}
    \item Jelaskan mengapa segmentasi diperlukan pada 8086. Kaitkan
          jawaban dengan lebar register segmen, lebar bus alamat,
          dan besar memori yang ingin dijangkau.

    \item Hitung alamat fisik dari pasangan berikut dan tuliskan
          setiap langkahnya.
          \begin{enumerate}
              \item \code{0000h:00A0h}
              \item \code{2000h:0010h}
              \item \code{9FFFh:0010h}
              \item \code{B800h:0080h}
          \end{enumerate}

    \item Jelaskan pengertian \eng{overlapping segmen}. Hitung
          banyaknya pasangan yang menunjuk alamat fisik
          \hex{00040}, tuliskan seluruh pasangan tersebut, dan
          jelaskan mengapa jumlahnya tidak pernah melebihi 4.096.

    \item Gambarkan pembagian memori 1 MB dan jelaskan kegunaan
          empat area: memori konvensional, buffer video, ROM
          adaptor, dan ROM BIOS. Jelaskan pula mengapa alamat
          \hex{B8000} termasuk area video.
\end{enumerate}

\subsection{Kisi-kisi Penilaian}

\begin{table}[H]
    \centering
    \caption{Kisi-kisi evaluasi Bab~\ref{chap:segmentasi}.}
    \label{tab:kisi-bab6}
    \small
    \begin{tabular}{@{}p{1.5cm}p{5.0cm}p{3.0cm}c@{}}
        \toprule
        \textbf{Sub-CPMK} & \textbf{Indikator yang diukur} &
        \textbf{Bentuk soal} & \textbf{Bobot} \\
        \midrule
        1.4 & Menjelaskan alasan segmentasi & Uraian 1 & 20\% \\
        \addlinespace
        1.4 & Menghitung alamat fisik &
        Pilihan ganda 1--3; Uraian 2 & 35\% \\
        \addlinespace
        1.4 & Menjelaskan paragraf dan \eng{overlapping} &
        Pilihan ganda 4, 5; Uraian 3 & 25\% \\
        \addlinespace
        1.4 & Menjelaskan pembagian memori 1 MB & Uraian 4 & 20\% \\
        \bottomrule
    \end{tabular}
\end{table}

\begin{checklistbox}
\vspace{1.1em}
\textbf{Checklist Pencapaian Kompetensi Bab~\ref{chap:segmentasi}}

Seorang mahasiswa dinyatakan kompeten pada Sub-CPMK 1.4 apabila:
\begin{itemize}
    \item[\checked] dapat menjelaskan mengapa 8086 memakai pasangan
          \eng{segment:offset};
    \item[\checked] dapat menghitung alamat fisik dengan menuliskan
          langkahnya;
    \item[\checked] dapat menjelaskan paragraf 16 byte dan segmen
          64 KB;
    \item[\checked] dapat menghitung banyaknya pasangan yang
          menunjuk satu alamat fisik;
    \item[\checked] dapat menyebutkan pembagian memori 1 MB;
    \item[\checked] dapat memverifikasi alamat dengan DEBUG atau
          emu8086.
\end{itemize}
\end{checklistbox}


\end{document}
